\chapter{Discussion}

\section{Model Capacity and Frequency Structure}

In the baseline scenario, almost all models achieve very high predictive
performance. The performance increase of the VQC when moving from
$L=1$ to $L=3$ (unary) or $L=2$ (ternary), followed by only marginal
improvements thereafter, can be explained by the previously analyzed
frequency structure. The target functions are dominated by
low-frequency components, so that even small encoding depths capture
most of the relevant structure. Increasing the depth further expands
the representable frequency spectrum, but provides only limited
additional benefit since high-frequency components are only weakly
present in the target functions.

The lower performance of the ternary models can also be explained by
their smaller number of effective parameters. Since ternary encoding
can represent a comparable frequency spectrum with fewer layers, the
number of effectively usable parameters is more strongly limited by
the achievable DLA.

Systematic differences can also be observed between the different
payoff structures. The Worst-of payoff represents the most challenging
approximation task. In contrast to Best-of and Average payoffs, the
function is not only bounded below by zero but is additionally defined
by a minimum operator. This structure introduces stronger nonlinearities
and therefore leads to slightly lower median $R^2$ values. These
nonlinearities are particularly pronounced near the payoff boundary,
where the option value transitions from zero to positive values.
Outside this region, the pricing function often behaves approximately
linearly, since changes in the input variables affect the payoff more
proportionally.

In the Worst-of scenario, the VQC also exhibits systematic positive
errors at low relative strikes in the upper value range, corresponding
to an underestimation of high option prices. In this region, the payoff
depends strongly on the minimum of the underlying assets, so that small
changes in the inputs can lead to larger changes in the function value.

Rotation-based VQCs can be interpreted functionally as models with a
finite Fourier-like frequency spectrum. Non-smooth or kink-like
structures can therefore only be approximated, while the globally
coupled parametrization induces a certain smoothing effect. This is
consistent with the observed underestimation of high target values
(Figure~\ref{fig:error-analysis}).

\section{Impact of Model Capacity under Limited Data}

With reduced training data, a more differentiated picture emerges.
Large Ferguson architectures maintain high predictive performance even
for very small datasets, suggesting that their larger model capacity
already provides an advantage with limited data and does not lead to
pronounced overfitting effects.

The Culkin model also shows stable performance for small datasets,
which is likely related to the regularizing effect of the applied
dropout.

In contrast, the VQC models exhibit an opposite pattern. The shallow
configuration ($L=1$) is the most robust for very small datasets,
whereas deeper variants show noticeable performance degradation. This
can be explained by the strongly expanded frequency representation at
larger encoding depths. If the training data are insufficient to
reliably determine additional degrees of freedom, unstable
approximations may occur. The shallow configuration therefore acts
like a more strongly regularized model.

As the amount of training data increases, this picture changes.
Once sufficient data are available, the deeper VQC architectures also
benefit from their larger model capacity and achieve predictive
performance comparable to classical networks.

\section{Robustness to Noise}

The experiments with artificial noise show that models with high
capacity achieve the best results in the noise-free case, but suffer
disproportionate performance losses as noise intensity increases. This
applies both to large Ferguson architectures and to VQCs with greater
encoding depth.

Nevertheless, the classical neural networks remain relatively robust
to moderate noise and still achieve high test-$R^2$ values even at
$\alpha = 0.5$.

The Culkin model with dropout shows particularly high stability,
which can plausibly be attributed to the regularizing effect of
dropout.

The shallow VQC ($L=1$) also exhibits only minor performance losses,
whereas deeper variants are considerably more sensitive to noisy
training data.

These observations can be explained by the effective model complexity.
Models with greater capacity possess more degrees of freedom and can
therefore more easily adapt to noise in the training data, whereas
simpler architectures act as stronger regularizers and thus remain
more robust.

\section{Interaction of Data Scarcity and Noise}

In the combined scenario of limited training data and strong noise,
this pattern becomes even more apparent. Models with lower capacity
achieve more stable results for very small datasets than more complex
architectures.

In particular, the shallow VQC remains one of the most robust
configurations across all considered sample sizes, whereas deeper
variants perform substantially worse when the dataset is very small.

As the amount of training data increases, however, this picture
changes considerably. Once sufficient data are available, the deeper
VQC models benefit from their greater expressive power and strongly
improve their predictive performance.

For the classical architectures, a similar but not identical effect
can be observed. While smaller Ferguson models benefit substantially
from increasing dataset size, very large models show only limited
improvement under strong noise when additional training data are
provided.


\section{Generalization under a Temporally Ordered Split}

Such a split is also closer to the real application scenario in
financial contexts, since models are typically trained on historical
data and applied to future market periods. While the classical
architectures largely retain their high predictive performance,
the VQC models show increased sensitivity, particularly in the
Best-of scenario.

The analysis shows that higher volatility levels occur during the
test period than were observed during training. The largest prediction
errors concentrate in these phases, indicating a distribution shift
(out-of-distribution) in the volatility features. In the Best-of case,
target values also appear that systematically exceed the range
observed during training.

This effect is particularly pronounced in the Best-of scenario.
Higher volatility increases the probability of extreme terminal
values of individual underlyings. Since the Best-of payoff is
determined by the maximum, this leads to more pronounced extreme
values and a steeper functional behavior than for Worst-of or
Average payoffs. Consistent with this, the performance degradation
of the VQC models is considerably smaller for Worst-of.

The deterioration is particularly visible for deeper VQC
configurations. One possible explanation is that greater encoding
depth improves interpolation within the training distribution,
but at the same time increases sensitivity to distribution shifts.
Classical ReLU-based architectures approximate functions in a
piecewise linear manner and extend learned slopes beyond the
training region, which may lead to more robust generalization in
this scenario.

\section{Limitations}

The original goal of this work was to study basket option prices
based on real observed market data, in order to represent the
correlations between the underlying assets using entanglement.
However, since hardly any freely available historical datasets
exist for actual basket option prices, the target values were
generated synthetically using Monte Carlo simulations.

Although the simulation is based on real historical time series
of the underlying assets, it represents a model-based
approximation. In particular, simplifying assumptions such as
constant correlations and simplified dynamics are made.
Consequently, the dependencies between the assets represented
through entanglement may differ from those observed in real
market prices.

The resulting target values therefore reflect the complexity of
real market prices only to a limited extent.

\section{Future Work}

A natural extension of this work is the use of more complex and
realistic datasets. In the present experiments, basket option
prices were generated using a simplified simulation model whose
frequency structure is dominated by low-frequency components.
Future work could investigate how the models behave when the
target functions exhibit more pronounced high-frequency
structures.

One possible direction is the use of real market data for basket
options, provided that suitable datasets become available.
Alternatively, more realistic simulation models could be
employed, for example incorporating time-dependent correlations,
stochastic volatility, or varying maturities.

Furthermore, it would be interesting to analyze to what extent
modeling absolute basket prices instead of relative returns
changes the functional structure of the target variable and
therefore imposes different requirements on the frequency
representation.

Another extension is the consideration of a larger number of
underlying assets. As dimensionality increases, the structural
complexity of the pricing function grows substantially. Since
the representable frequency spectrum of rotation-based VQCs
increases exponentially with the input dimension, more
pronounced differences between classical and quantum-based
models may emerge in higher-dimensional settings.

Finally, it could be investigated whether appropriate
preprocessing of the target variable improves modeling.
Since option prices are nearly linear across large regions of
the input space and stronger nonlinearities mainly occur near
the payoff boundaries, a transformation such as an arccos
mapping could help distribute the functional structure more
evenly and thereby facilitate approximation by rotation-based
VQCs.